% ---------------------------------------------------------------------------
% Documentação técnica do quad-site.
%
% Documento autocontido: compila com pdflatex sem pacotes fora do TeX Live
% básico. Para embutir em um relatório maior, remova o preâmbulo e o
% \begin{document}/\end{document}, mantendo apenas as seções.
% ---------------------------------------------------------------------------
\documentclass[11pt,a4paper]{article}

\usepackage[utf8]{inputenc}
\usepackage[T1]{fontenc}
\usepackage[brazil]{babel}
\usepackage[margin=2.5cm]{geometry}
\usepackage{booktabs}
\usepackage{longtable}
\usepackage{array}
\usepackage{xcolor}
\usepackage{listings}
\usepackage[hidelinks]{hyperref}

\makeatletter
\newcommand{\cd@us}{\textunderscore\discretionary{}{}{}}
\newcommand{\cd}[1]{{\ttfamily\hyphenchar\font\m@ne\let\_\cd@us #1}}
\makeatother

\definecolor{codebg}{gray}{0.96}
\definecolor{codecomment}{rgb}{0.35,0.45,0.35}
\definecolor{codekey}{rgb}{0.20,0.20,0.60}

\lstdefinestyle{base}{
  backgroundcolor=\color{codebg},
  basicstyle=\ttfamily\footnotesize,
  breaklines=true,
  breakatwhitespace=true,
  columns=fullflexible,
  keepspaces=true,
  frame=single,
  framerule=0pt,
  xleftmargin=6pt,
  xrightmargin=6pt,
  aboveskip=8pt,
  belowskip=8pt,
  showstringspaces=false,
  commentstyle=\color{codecomment},
  keywordstyle=\color{codekey}\bfseries,
  extendedchars=true,
  inputencoding=utf8,
  literate=
    {á}{{\'a}}1 {é}{{\'e}}1 {í}{{\'i}}1 {ó}{{\'o}}1 {ú}{{\'u}}1
    {Á}{{\'A}}1 {É}{{\'E}}1 {Í}{{\'I}}1 {Ó}{{\'O}}1 {Ú}{{\'U}}1
    {à}{{\`a}}1 {À}{{\`A}}1
    {â}{{\^a}}1 {ê}{{\^e}}1 {ô}{{\^o}}1
    {Â}{{\^A}}1 {Ê}{{\^E}}1 {Ô}{{\^O}}1
    {ã}{{\~a}}1 {õ}{{\~o}}1 {Ã}{{\~A}}1 {Õ}{{\~O}}1
    {ç}{{\c c}}1 {Ç}{{\c C}}1
    {ü}{{\"u}}1 {Ü}{{\"U}}1,
}
\lstdefinestyle{tsx}{style=base,language=}
\lstdefinestyle{shell}{style=base,language=bash}
\lstdefinestyle{plain}{style=base,language=}

\title{Documentação técnica do \cd{quad-site}\\
       \large Quad Concursos --- site institucional}
\author{}
\date{}

\begin{document}
\maketitle

\tableofcontents
\newpage

% ===========================================================================
\section{Objetivo e escopo}

Este documento descreve o estado atual do \cd{quad-site}, o site institucional
da Quad Concursos (Feira de Santana, Bahia): sua arquitetura, o sistema de
design que rege a interface, as convenções de componentes e dados, a forma
como mudanças são verificadas neste projeto, e um guia de como formular
solicitações a um agente de IA para obter resultados consistentes com o que já
existe.

O objetivo é permitir que qualquer pessoa --- humana ou agente --- retome o
trabalho neste código sem precisar redescobrir por tentativa e erro decisões já
tomadas: por que as cores são as que são, por que não há testes automatizados,
por que o checkout não está ligado ao carrinho, e o que fazer quando uma
mudança visual precisa ser verificada sem acesso a um navegador.

A seção \ref{sec:limitacoes} lista de forma explícita o que \textbf{não} existe
neste projeto, para que o leitor não presuma infraestrutura (testes, CI,
backend) que nunca foi construída.

% ===========================================================================
\section{Visão geral do projeto}
\label{sec:visao-geral}

\subsection{Stack}

\begin{longtable}{@{}p{4.0cm}p{10.5cm}@{}}
\toprule
\textbf{Camada} & \textbf{Tecnologia} \\
\midrule
\endhead
Framework & Next.js 16 (App Router), React 19 \\
Linguagem & TypeScript \\
Estilo & Tailwind CSS v3, sem biblioteca de componentes \\
Fontes & Rajdhani (\cd{--font-display}) e Work Sans (\cd{--font-body}), via
\cd{next/font/google} ou \cd{<link>} no \cd{layout.tsx} \\
Estado do carrinho & Contexto React (\cd{lib/cart.tsx}) persistido em
\cd{localStorage} \\
Backend & Nenhum \\
Hospedagem prevista & Não definida no repositório \\
\bottomrule
\end{longtable}

\subsection{Decisão deliberada: sem backend por enquanto}

Todo o conteúdo do site --- turmas, produtos, eventos, professores, aprovados,
posts do blog --- vem de arrays tipados estáticos em \cd{lib/data.ts}. Esta é
uma decisão de sequenciamento, não uma limitação técnica: o frontend é
finalizado primeiro; a persistência real (banco de dados, upload de imagens em
bucket) entra depois, sem que a modelagem de dados precise mudar --- os campos
opcionais como \cd{fotoUrl} já existem nos tipos, prontos para receber uma URL
real quando a infraestrutura existir.

Consequência prática para quem for integrar um backend futuramente: trocar a
origem dos arrays de \cd{lib/data.ts} por chamadas a uma API mantém os
componentes que os consomem inalterados, pois eles já operam sobre os tipos
exportados (\cd{Turma}, \cd{Produto}, \cd{Evento} etc.), não sobre os arrays em
si.

% ===========================================================================
\section{Arquitetura}

\subsection{Estrutura de pastas}

\begin{longtable}{@{}p{3.2cm}p{11.3cm}@{}}
\toprule
\textbf{Pasta} & \textbf{Conteúdo} \\
\midrule
\endhead
\cd{app/} & Rotas do App Router. Cada \cd{page.tsx} é uma rota; pastas entre
colchetes (\cd{[slug]}) são rotas dinâmicas. \\
\cd{components/} & Componentes de UI compartilhados entre rotas. Sem
subpastas: a convenção é um arquivo por componente, nome em PascalCase igual
ao componente exportado. \\
\cd{lib/} & \cd{data.ts} (dados mockados e seus tipos), \cd{cart.tsx}
(contexto do carrinho), \cd{format.ts} (utilitários de formatação
compartilhados). \\
\cd{public/images/} & Ativos estáticos servidos na raiz (\cd{/images/...}):
logo em SVG, fotos de professores em WebP/PNG. \\
\bottomrule
\end{longtable}

\subsection{Rotas existentes}

\begin{longtable}{@{}p{5.4cm}p{9.1cm}@{}}
\toprule
\textbf{Rota} & \textbf{Conteúdo} \\
\midrule
\endhead
\cd{/} & Home: hero, turmas, PDFs, eventos, carrossel de professores,
aprovados, blog \\
\cd{/turmas}, \cd{/produtos}, \cd{/eventos} & Listagens \\
\cd{/turmas/[slug]}, \cd{/produtos/[slug]}, \cd{/eventos/[slug]} & Detalhe de
item, com \cd{BuyBox} \\
\cd{/professores/[slug]} & Banner em modo paisagem com foto do professor,
credenciais \\
\cd{/aprovados} & Cards de aprovados e depoimentos em vídeo \\
\cd{/blog}, \cd{/blog/[slug]} & Listagem e post \\
\cd{/cursos} & Cursos online (redireciona para a Tutory) \\
\cd{/carrinho}, \cd{/checkout}, \cd{/checkout/confirmacao} & Fluxo de compra \\
\cd{/entrar}, \cd{/cadastro} & Login (mockado) e cadastro em duas etapas \\
\cd{/aluno} e subrotas & Área do aluno: painel, materiais, turmas, eventos,
compras, perfil \\
\bottomrule
\end{longtable}

\subsection{Modelo de dados}

\cd{lib/data.ts} exporta um tipo e um array por entidade: \cd{Turma},
\cd{Produto}, \cd{Evento}, \cd{Aprovado}, \cd{Professor}, \cd{Post},
\cd{HistoriaVideo}. Cada tipo carrega apenas os campos que a interface
efetivamente usa hoje --- não há campos especulativos para uma versão futura.

\begin{lstlisting}[style=tsx]
export type Professor = {
  slug: string;
  nome: string;
  materia: string;
  bio: string;
  bioCompleta: string;
  formacao: string[];
  fotoUrl?: string; // opcional: sem foto real, cai no fallback de iniciais
};
\end{lstlisting}

O padrão \cd{fotoUrl?: string} se repete em \cd{Professor} e \cd{Aprovado}: o
componente que consome o campo sempre trata a ausência de foto renderizando um
monograma com as iniciais (\cd{lib/format.ts}, função \cd{getIniciais}) sobre
um gradiente, em vez de quebrar ou deixar um espaço vazio.

\subsection{Estado do carrinho}

O carrinho é um \cd{CartProvider} (Contexto React) que envolve toda a árvore em
\cd{layout.tsx}, com persistência em \cd{localStorage} sob a chave
\cd{quad-carrinho}. Ele é agnóstico ao tipo de item --- produtos, turmas e
eventos compartilham a mesma estrutura \cd{CartItem}, com uma regra de negócio
única: apenas itens do tipo \cd{"produto"} incrementam quantidade ao serem
adicionados de novo; turmas e eventos são de item único.

\begin{lstlisting}[style=tsx]
function addItem(item: Omit<CartItem, "quantidade">) {
  setItems((prev) => {
    const existente = prev.find((i) => i.id === item.id);
    if (existente) {
      if (item.tipo !== "produto") return prev; // turma/evento: item único
      return prev.map((i) => (i.id === item.id ? { ...i, quantidade: i.quantidade + 1 } : i));
    }
    return [...prev, { ...item, quantidade: 1 }];
  });
}
\end{lstlisting}

O checkout, deliberadamente, \textbf{não} lê esse contexto: o resumo do pedido
em \cd{/checkout} é hardcoded. Ver seção \ref{sec:limitacoes}.

% ===========================================================================
\section{Sistema de design}
\label{sec:design}

Esta é a base de conhecimento do projeto: em vez de regras em Markdown, ela vive
como tokens no \cd{tailwind.config.ts} e classes utilitárias em
\cd{globals.css}. Editar um token aqui propaga a mudança para todo lugar que o
usa --- é o mesmo princípio de fonte única de verdade de uma base de regras
versionada, só que expresso em configuração de build em vez de texto.

\subsection{Paleta de cores}

\begin{longtable}{@{}p{3.4cm}p{3.0cm}p{7.7cm}@{}}
\toprule
\textbf{Token} & \textbf{Valor} & \textbf{Uso} \\
\midrule
\endhead
\cd{navy} / \cd{navy-deep} / \cd{navy-mid} &
\#0D1B2A / \#091320 / \#14263A & Fundos escuros, do mais claro ao mais escuro \\
\cd{blue} / \cd{blue-dark} & \#4C8DFF / \#2F63C4 & Cor de destaque (CTAs,
links, badges) \\
\cd{offwhite} & \#F5F3EE & Fundo padrão do corpo da página \\
\cd{steel} / \cd{steel-light} & \#8B95A3 / \#B8C0CC & Texto secundário, em
fundo claro e escuro respectivamente \\
\cd{border-light} / \cd{border-dark} & \#E2DFD6 / \cd{rgba(255,255,255,.1)} &
Bordas em fundo claro e escuro \\
\bottomrule
\end{longtable}

Esta paleta é considerada estável: várias rodadas de ajuste visual nesta sessão
mantiveram essas cores intactas por instrução explícita do responsável pelo
produto.

\subsection{Escalas de espaçamento, tipografia e elevação}

\begin{longtable}{@{}p{3.0cm}p{4.0cm}p{7.1cm}@{}}
\toprule
\textbf{Escala} & \textbf{Tokens} & \textbf{Papel} \\
\midrule
\endhead
Espaçamento vertical & \cd{py-header} (56px), \cd{py-section} (72px) &
Substituem valores soltos (\cd{py-12}, \cd{py-16}, \cd{py-[72px]}) que
existiam antes da padronização \\
Largura de container & \cd{max-w-site} (1120px), \cd{max-w-narrow} (760px),
\cd{max-w-form} (920px) & Largura padrão, leitura de texto corrido (posts,
corpo de turma/produto), formulários transacionais (checkout) \\
Tipografia de título & \cd{text-display}, \cd{text-h1}, \cd{text-h1-sm},
\cd{text-h2} & Hero da home, heróis de entidade, páginas utilitárias,
\cd{SectionHead} \\
Elevação & \cd{shadow-sm}, \cd{shadow} (padrão), \cd{shadow-md} & Aplicadas
apenas a superfícies que sobrepõem conteúdo (mega menu, \cd{BuyBox} fixo,
overlay de hover); cards estáticos permanecem sem sombra, só com borda \\
\bottomrule
\end{longtable}

\subsection{Classes utilitárias em \cd{globals.css}}

\begin{longtable}{@{}p{3.2cm}p{11.3cm}@{}}
\toprule
\textbf{Classe} & \textbf{Comportamento} \\
\midrule
\endhead
\cd{.btn-blue} / \cd{.btn-outline} / \cd{.btn-outline-dark} &
Botão com efeito de profundidade física: sombra deslocada que recua ao clicar
(\cd{box-shadow} + \cd{translateY} no \cd{:hover}/\cd{:active}). É a
assinatura de interação do site; não deve ser redesenhada sem pedido
explícito. \\
\cd{.eyebrow} & Rótulo azul, maiúsculo, com um tracinho antes do texto.
Reutilizado tanto como categoria de página quanto como cabeçalho de
subseção dentro de uma mesma página (ex.\ a página de professor usa
\cd{.eyebrow} duas vezes: matéria do professor e ``Formação e experiência''). \\
\cd{.hexagon} & \cd{clip-path} hexagonal, usado em avatares, marcadores de
lista (\cd{CheckList}) e no ícone de confirmação do checkout. \\
\cd{.focus-ring} & \cd{outline} azul em \cd{:focus-visible}, para todo
elemento interativo que não seja \cd{.btn} (que já tem a própria regra de
foco). \\
\cd{.pattern-diagonal} & Textura de listras diagonais sutis, usada no hero da
home e nos placeholders da \cd{Gallery}. \\
\bottomrule
\end{longtable}

\subsection{Convenções de componente}

\begin{itemize}
  \item \textbf{Prop \cd{dark}}: vários componentes (\cd{AccentCard},
        \cd{ProductCard}, \cd{SectionHead}, \cd{CheckList}) aceitam uma prop
        booleana \cd{dark} que alterna a paleta de texto/borda para uso sobre
        fundo \cd{navy}. É o mecanismo padrão de reuso entre seções claras e
        escuras --- criar uma variante nova do componente em vez de estender
        essa prop é considerado duplicação.
  \item \textbf{\cd{Reveal}} (\cd{components/Reveal.tsx}): wrapper de entrada
        suave via \cd{IntersectionObserver}, respeitando
        \cd{prefers-reduced-motion}. Aplicado a cards abaixo da dobra, nunca a
        conteúdo acima da dobra, sticky, ou páginas transacionais
        (\cd{/carrinho}, \cd{/checkout}, \cd{/entrar}, \cd{/cadastro}).
  \item \textbf{Máscaras de campo}: não há biblioteca de máscara de input.
        CPF, telefone e CEP (\cd{app/cadastro/page.tsx}) usam funções puras
        que limpam não-dígitos e reconstroem a pontuação a cada tecla,
        aplicadas via uma prop \cd{mask?: (v: string) => string} no
        componente \cd{Field} local dessa página.
  \item \textbf{Nomenclatura}: identificadores de domínio (variáveis, props,
        dados) são em português (\cd{nome}, \cd{preco}, \cd{tipo},
        \cd{criarConta}); nomes de componentes e conceitos técnicos
        genéricos são em inglês (\cd{Field}, \cd{Reveal}, \cd{dark}). Essa
        mistura é intencional e consistente em todo o código existente.
\end{itemize}

% ===========================================================================
\section{Verificação}
\label{sec:verificacao}

Este projeto \textbf{não possui} suíte de testes automatizados, pipeline de
CI, nem hooks de Git ou de agente configurados. A verificação de uma mudança é
inteiramente manual, e se apoia em três mecanismos:

\begin{enumerate}
  \item \textbf{Checagem de tipos}: \cd{npx tsc --noEmit}. É o único portão
        automatizado real do projeto --- captura classes Tailwind referenciando
        props inexistentes, imports quebrados e incompatibilidades de tipo.
        Deve ser executado depois de qualquer edição em arquivo \cd{.tsx}.
  \item \textbf{Checagem de rota}: com o servidor de desenvolvimento
        (\cd{npm run dev}) já em execução, \cd{curl} contra cada rota tocada,
        conferindo código \cd{200} e, quando relevante, que um trecho de texto
        esperado apareça no HTML renderizado no servidor.
  \item \textbf{Inspeção do CSS compilado}: quando uma classe Tailwind nova é
        introduzida via token (\cd{padding}, \cd{maxWidth}, \cd{fontSize},
        \cd{boxShadow} customizados), buscar a regra correspondente no arquivo
        CSS servido em \cd{/\_next/static/chunks/...css} confirma que o
        Tailwind realmente gerou a classe, em vez de presumir que o nome do
        token está correto.
\end{enumerate}

\begin{lstlisting}[style=shell]
npx tsc --noEmit
curl -s -o /dev/null -w "%{http_code} /rota\n" http://localhost:3000/rota
\end{lstlisting}

\subsection{\cd{next lint} está quebrado neste ambiente}

O script \cd{npm run lint} (\cd{next lint}) falha com
\texttt{"Invalid project directory provided, no such directory: ...\textbackslash lint"}
neste ambiente Windows, independentemente do código --- é um problema de
parsing de argumento do CLI do Next.js com o caminho do projeto contendo
espaços, não algo introduzido por uma mudança específica. Na ausência de lint
funcional, \cd{tsc --noEmit} é o único checador estático disponível.

\subsection{Sem automação de navegador}

Não há Playwright, Puppeteer, nem qualquer ferramenta de captura de tela neste
fluxo de trabalho. Mudanças puramente visuais (posição de elemento, corte de
imagem, animação, hover) não podem ser verificadas de forma independente ---
apenas o HTML/CSS servido pode ser inspecionado via \cd{curl}, o que confirma
que uma classe foi aplicada, mas não como ela renderiza. A verificação visual
final depende de captura de tela fornecida por quem está pedindo a mudança.

% ===========================================================================
\section{Fluxo de trabalho}

Diferente de um pipeline com portões automatizados, o ciclo de trabalho neste
projeto é conversacional e iterativo:

\begin{longtable}{@{}p{4.6cm}p{10.0cm}@{}}
\toprule
\textbf{Momento} & \textbf{Comportamento observado} \\
\midrule
\endhead
Pedido inicial & Frequentemente descreve o resultado desejado, não a
implementação. Ajustes finos de posicionamento/corte costumam vir em rodadas
subsequentes, guiados por captura de tela. \\
\addlinespace
Edição concorrente & Quem pede a mudança pode editar o mesmo arquivo
diretamente no editor enquanto o agente trabalha. O agente deve reler o
arquivo antes de cada edição em vez de assumir que seu último estado
conhecido ainda é válido. \\
\addlinespace
Após cada edição relevante & \cd{tsc --noEmit} e, quando aplicável,
\cd{curl} na rota afetada (seção \ref{sec:verificacao}). \\
\addlinespace
Ajuste visual fino & Quando um ajuste de pixel específico (posição, corte de
imagem, zoom) não converge em uma ou duas tentativas, vale expor o mecanismo
subjacente (ex.\ largura da caixa controla zoom via \cd{object-cover}) para
que a calibração final possa ser feita ao vivo no DevTools do navegador, sem
depender de rodadas adicionais. \\
\bottomrule
\end{longtable}

Não há \cd{commit} automático: alterações ficam no diretório de trabalho até
serem revisadas e commitadas explicitamente.

% ===========================================================================
\section{Guia de formulação de solicitações}
\label{sec:prompts}

\subsection{Princípio}

Assim como em qualquer projeto com convenções já estabelecidas, reenunciar o
sistema de design a cada pedido (``use a cor azul de destaque'', ``siga o
padrão de card com borda esquerda'') é redundante --- essas decisões já estão
no código existente e serão seguidas por consistência com o que já existe. O
que vale a pena declarar é o que \textbf{não} pode ser inferido do código: o
conteúdo específico, a hierarquia da informação, e --- principalmente neste
projeto, que não tem testes --- o critério de ``pronto'' para algo visual.

\subsection{Nova entidade de conteúdo (seguindo o padrão \cd{lib/data.ts})}

\begin{lstlisting}[style=plain]
Adicione um novo tipo Depoimento em lib/data.ts, seguindo o mesmo padrão
de Aprovado (tipo exportado + array exportado).

Campos: nome, cargo, texto (até 280 caracteres), fotoUrl opcional.

Crie um componente DepoimentoCard reaproveitando o padrão de fallback de
iniciais já usado em AprovadoCard/ProfessoresCarousel (getIniciais de
lib/format.ts) quando fotoUrl não existir.

Adicione uma seção "Depoimentos" na home, entre Aprovados e Blog, com o
mesmo ritmo de espaçamento (py-section) das seções vizinhas.
\end{lstlisting}

Note o que o pedido \emph{não} especifica: cor, borda, sombra, ou onde a
seção deve ficar no arquivo \cd{tailwind.config.ts} --- isso já está
resolvido pelo sistema de design existente. O que ele especifica é o que é
específico do domínio: os campos, o limite de caracteres, e a posição na
página.

\subsection{Ajuste visual fino}

Pedidos de posicionamento/recorte de imagem se beneficiam de declarar o
resultado desejado e, quando possível, anexar uma captura de tela --- dado que
não há automação de navegador (seção \ref{sec:verificacao}), a captura é o
único canal de verificação visual real.

\begin{lstlisting}[style=plain]
[captura de tela anexada]

A cabeça do professor está cortada no topo do banner. Quero ver o rosto
inteiro, mesmo que isso corte mais do ombro/peito na base.
\end{lstlisting}

\subsection{Mudança de comportamento existente}

Ao alterar um fluxo já implementado, enunciar o estado atual reduz ambiguidade
--- sobretudo porque não há teste automatizado que documente o comportamento
anterior:

\begin{lstlisting}[style=plain]
Hoje, ao adicionar um produto já presente no carrinho, a quantidade é
incrementada; para turma/evento, adicionar de novo não faz nada (item
único). Mantenha esse comportamento, mas mostre um aviso rápido
("Já está no carrinho") quando o clique em turma/evento for ignorado.
\end{lstlisting}

\subsection{Formulações a evitar}

\begin{longtable}{@{}p{6.6cm}p{7.9cm}@{}}
\toprule
\textbf{Formulação} & \textbf{Problema} \\
\midrule
\endhead
``Deixe essa página mais bonita'' &
Sem critério verificável. Nenhum teste captura ``bonita''; o pedido deveria
nomear o elemento específico e o resultado esperado. \\
\addlinespace
``Adicione testes para isso'' &
Não há framework de teste configurado neste projeto (Jest, Vitest, Playwright
--- nenhum instalado). O pedido precisa primeiro decidir se vale introduzir
essa infraestrutura, o que é uma mudança de escopo maior, não um ajuste
pontual. \\
\addlinespace
``Conecte isso a um banco de dados'' &
Contraria a decisão deliberada de manter o frontend mockado até estar
completo (seção \ref{sec:visao-geral}). Requer confirmação explícita antes de
proceder. \\
\addlinespace
``Use uma biblioteca de máscara de input para o telefone'' &
O padrão do projeto é função pura sem dependência (seção
\ref{sec:design}); uma biblioteca nova deve ser uma decisão deliberada, não
um atalho de conveniência. \\
\bottomrule
\end{longtable}

% ===========================================================================
\section{Limitações conhecidas}
\label{sec:limitacoes}

\begin{itemize}
  \item \textbf{Sem backend nem banco de dados.} Todo conteúdo é mockado em
        \cd{lib/data.ts}. Decisão deliberada, não pendência esquecida (ver
        seção \ref{sec:visao-geral}).
  \item \textbf{Sem armazenamento de imagens em bucket.} Fotos de professores
        ficam em \cd{public/images/}, versionadas junto ao código. Migração
        futura recomendada para um bucket (Vercel Blob ou Cloudinary, dado o
        deploy provável em Vercel) mantendo apenas a URL no dado, não o
        arquivo em si.
  \item \textbf{Sem suíte de testes automatizados.} Nenhum framework de teste
        instalado. Toda verificação é manual (seção \ref{sec:verificacao}).
  \item \textbf{Sem CI, sem hooks de Git, sem hooks de agente.} Não existe
        portão que bloqueie um commit ou um push por violação de padrão --- a
        consistência depende inteiramente de quem edita o código seguir as
        convenções documentadas aqui.
  \item \textbf{\cd{next lint} não funciona neste ambiente} (seção
        \ref{sec:verificacao}), então nenhuma checagem de estilo ESLint roda
        atualmente, automatizada ou manual.
  \item \textbf{Checkout não lê o carrinho real.} O resumo em \cd{/checkout}
        é hardcoded; decisão explícita para manter o escopo do carrinho
        contido até uma integração de pagamento real ser definida.
  \item \textbf{Sem automação de navegador.} Nenhuma mudança puramente visual
        pode ser verificada de forma independente; depende de captura de tela
        de quem solicita.
\end{itemize}

% ===========================================================================
\section{Estado atual}

Levantamento no momento da escrita deste documento:

\begin{longtable}{@{}p{5.0cm}p{9.5cm}@{}}
\toprule
\textbf{Métrica} & \textbf{Valor} \\
\midrule
\endhead
Rotas & 23 arquivos \cd{page.tsx}, incluindo 5 rotas dinâmicas
(\cd{[slug]}) \\
Componentes compartilhados & 21, em \cd{components/} \\
Tipos de entidade & 7 (\cd{Turma}, \cd{Produto}, \cd{Evento}, \cd{Aprovado},
\cd{HistoriaVideo}, \cd{Professor}, \cd{Post}) \\
Dependências de produção & 3: \cd{next}, \cd{react}, \cd{react-dom} ---
nenhuma biblioteca de UI, formulário, máscara ou requisição HTTP além do
\cd{fetch} nativo \\
Integrações externas & ViaCEP (busca de endereço por CEP em
\cd{/cadastro}), sem chave de API \\
\bottomrule
\end{longtable}

Este levantamento não substitui a leitura do código: é um retrato do momento
da escrita, sujeito a ficar desatualizado conforme o projeto evolui. Antes de
confiar em qualquer número acima para uma decisão, confirme contra o estado
atual do repositório.

\end{document}
